\documentclass[10pt,a4paper]{amsart}
\usepackage[margin=19mm]{geometry}
\usepackage{amsmath,amssymb,mathtools}
\usepackage[scr=rsfs,cal=euler]{mathalfa}
\usepackage{xfrac}
\usepackage[unicode,hidelinks,bookmarks=false]{hyperref}
\newcommand{\ie}{i.e.\ }
\newcommand{\cf}{cf.\ }
\newcommand{\resp}{respectively\ }
\newcommand{\Subsection}{\subsection}
\providecommand{\nobreakdash}{\nobreak\hbox{-}}
\setlength{\emergencystretch}{2em}
\multlinegap0pt
\newcommand{\RR}{\mathbb{R}}
\usepackage{amssymb}
\usepackage{tikz}
\usepackage{mathtools}
\usepackage{bm}
\usepackage{float}
\usepackage{enumerate}
\newtheorem{lemma}{Lemma}
\title{On colorings of hypergraphs embeddable in $\RR^d$}
\author{Seunghun Lee, Eran Nevo}
\date{Source: arXiv:2307.14195v3 (CC BY 4.0)}
\begin{document}
\maketitle

\noindent\emph{Source and licence.} On colorings of hypergraphs embeddable in $\RR^d$. Version arXiv:2307.14195v3 (CC BY 4.0), distributed by arXiv under the Creative Commons Attribution 4.0 International licence, http://creativecommons.org/licenses/by/4.0/, as recorded in the arXiv licence metadata for this exact version and shown on the abstract page https://arxiv.org/abs/2307.14195v3. Full licence notice: \texttt{licenses/arxiv-2307.14195-CC-BY-4.0.txt}.\par\smallskip

\noindent\textit{Editorial scope.} Counted unit: the article's lemma lemma\_base\_step\_H(a,b) (source0.tex lines 286--288). The article's complete original proof of it is delivered with it (source0.tex lines 289--306, one \texttt{proof} environment of 18 lines). No mathematics is written, repaired or re-derived: the statement and the proof are the article's own lines, in the article's own notation, and no retained line is substituted or re-flowed.  No further source-local context is needed: every cross-reference of every retained line resolves inside the two delivered spans.  Nothing was omitted from any retained span, so the package carries no omission record.  No retained line contains a citation, so the wrapper carries no bibliography and no bibliography entry.  No retained line contains a figure environment, an image-inclusion command or drawing code, so no image is delivered and the package declares no asset dependency.  Editorial formatting only, in addition: an AMS \texttt{amsart} preamble that restates the article's own declarations and macros used by the retained lines, namely the environments \texttt{lemma} on separate counters, together with the standard packages those lines need.  The retained environments print in this document as Lemma 1 in order of appearance, whereas the article prints the same items with the numbering of its own full document.  The complete native source of the article as distributed in the arXiv e-print for this exact version is bundled unchanged as \texttt{sources/144814/source0.tex}.
\begin{lemma} \label{lemma_base_step_H(a,b)}
  When $3\leq a \leq b$, $H(a,b)$ is not $(a+2)$-interlacing with respect to any DFS (depth first search) order on $T(a,b)$.
\end{lemma}

\begin{proof}
For $T=T(a,b)$ we fix a DFS order $\prec_T$. We choose two distinct hyperedges $A$ and $B$ and consider the following exhaustive three cases. 

\smallskip

\textbf{Case I.} $A$ and $B$ are both maximal chains: The order $\prec_T$ implies that the pattern of the chains $A$ and $B$ is either of the form $I^lA^{a-l}B^{a-l}$ or $I^lB^{a-l}A^{a-l}$ for some $1\leq l \leq a-1$ ($l\ge 1$ as both $A$ and $B$ contain the root of $T$). So $\{A, B\}$ is at most $(a+1)$-interlacing (when $l=a-1$), but cannot be $(a+2)$-interlacing.

\smallskip 

\textbf{Case II.} Exactly one of $A$ and $B$ is a maximal chain: We may assume that $A$ is a maximal chain. When $A$ and $B$ intersect, the order $\prec_T$ implies that the pattern of $A$ and $B$ is in the form of $A^lB^rIA^{a-l-1}B^{b-r-1}$ where $1\leq l \leq a$ and $0 \leq r \leq b$ ($l\geq 1$ as the root of $T$ cannot be in $B$). Then $\{A,B\}$ is at most 4-interlacing, thus $\{A,B\}$ is not $(a+2)$-interlacing as $a \geq 3$. When $A$ and $B$ are disjoint, the pattern of $A$ and $B$ with respect to $\prec_T$ is $A^lB^bA^{a-l}$ for $1\leq l \leq a$,
hence $\{A,B\}$ is at most 3-interlacing.

\smallskip
		
\textbf{Case III.} Neither $A$ nor $B$ is a chain: Let $A$ and $B$ be the set of children of non-leaf vertices $v$ and $w$, respectively. 
If $w$ is a descendant of $v$, then $\prec_T$ implies that the pattern of $A$ and $B$ is $A^lB^bA^{b-l}$  where $0\leq l \leq b$. 
If $v$ is a descendant of $w$ the pattern of $A$ and $B$ with respect to $\prec_T$ is $B^lA^bB^{b-l}$ where $0\leq l \leq b$. In both cases $\{A,B\}$ is just at most 3-interlacing. Else, $v$ and $w$ are not in the same maximal chain, then the pattern of $A$ and $B$ is either $A^bB^b$ or $B^bA^b$ , hence $\{A,B\}$ is just at most 2-interlacing. In either case, $\{A,B\}$ is not $(a+2)$-interlacing for $a\geq 3$.
\end{proof}

\end{document}
